% Isolated after CP34; no canonical or shared build input is changed.
\subsection{A repeated determinant in an actual twisted product pencil}

Work over an algebraically closed field of characteristic two. Let
\(E:y^2+y=x^3\), \(M_E=\mathcal O_E(3O)\), and
\(B=E\times E\), with \(M=M_E\boxtimes M_E\).
Let \(P=P_1\boxtimes\mathcal O_E\), where
\(P_1\in\operatorname{Pic}^0(E)\) is nontrivial. The following
concerns the actual relative Cartier cross pairing on
\(H^0(MP)\times H^0(MP^{-1})\). It does not replace relative
Cartier by a ground-field-linear absolute Cartier operator.

\begin{lemma}[The full nine-variable cross determinant is a cube]
\label{lem:ss-product-one-axis-cross-determinant-cube}
\status{Candidate; complete covariance proof and exact checks supplied.}
Use the target bases \((1,X,Y)\) and \((1,Z,W)\) of
Lemma~\ref{lem:DE-SS-cross-elliptic}. Write a target functional
as the \(3\times3\) coefficient matrix \(T=(t_{ij})\).
In the stated source bases, its nine-dimensional cross matrix
\(C(T)\) satisfies the exact polynomial identity
\[
                        \det C(T)=\det(T)^3.
 \tag{PTC1}
\]
This identity holds for every nontrivial \(P_1\), not just
for a finite list of twist parameters.
\end{lemma}
\begin{proof}
Write \(P_1=\mathcal O_E((a,b)-O)\), so \(b^2+b=a^3\).
The already proved relative Cartier calculation gives
\[
 A_0=\begin{pmatrix}1&a&0\\a&0&b\\0&b+1&a\end{pmatrix},
 \quad A_1=R=\begin{pmatrix}0&0&1\\0&1&0\\1&0&0\end{pmatrix},
 \quad A_2=\begin{pmatrix}0&0&0\\0&0&1\\0&1&0\end{pmatrix}.
 \tag{PTC2}
\]
On the untwisted second factor, the three coefficient matrices
\(B_0,B_1,B_2\) are the alternating matrices with their sole
nonzero upper-triangular entry respectively at \((1,2),(1,3),(2,3)\).
Thus they form the full space of alternating forms on a three-space,
and Kunneth gives
\[
                    C(T)=\sum_{i,j=0}^2t_{ij}A_i\otimes B_j.
 \tag{PTC3}
\]

For \(g\in\operatorname{GL}_3\), write
\(g^tB_jg=\sum_k S(g)_{kj}B_k\). Exterior-square covariance
gives \(\det S(g)=\det(g)^2\), and direct substitution in
(PTC3) gives
\[
 C(TS(g)^t)=(I_3\otimes g^t)C(T)(I_3\otimes g),
 \quad
 \det C(TS(g)^t)=\det S(g)^3\det C(T).
 \tag{PTC4}
\]
The map \(g\mapsto S(g)\) is surjective on geometric points.
Indeed, under the volume identification of the exterior square
with the dual three-space, it is conjugate to
\(g\mapsto\det(g)g^{-1}\). Given any invertible matrix \(h\),
choose \(\delta\) with \(\delta^2=\det h\), and take
\(g=\delta h^{-1}\) in these coordinates. This argument uses
algebraic closedness, not separability of the exterior-square map.

Right multiplication is transitive on invertible matrices \(T\),
so (PTC4) shows that \(\det C(T)/\det(T)^3\) is constant
there. At the permutation matrix whose nonzero entries are
\(t_{10}=t_{01}=t_{22}=1\), the cross matrix is, after ordering
the second factor outside,
\[
 \begin{pmatrix}0&R&A_0\\R&0&A_2\\A_0&A_2&0\end{pmatrix}.
\]
Its upper six-dimensional block is invertible, and its Schur
complement is
\[
 A_0RA_2+A_2RA_0
 =\begin{pmatrix}0&1&a\\1&0&0\\a&0&1\end{pmatrix},
 \qquad\det=1.
\]
This is the universal unit calculation of
Lemma~\ref{lem:DE-SS-cross-product}; the permutation matrix has
determinant one in characteristic two. The constant is therefore
one. Equality on the dense invertible locus proves (PTC1) as a
polynomial identity for every allowed \((a,b)\).
\end{proof}

\begin{proposition}[The determinant divisor has generic corank three]
\label{prop:ss-product-one-axis-cross-corank-three}
\status{Candidate; complete rank calculation supplied.}
If \(\operatorname{rank}T=2\), then \(\operatorname{rank}C(T)=6\).
Consequently the reduced determinant hypersurface is the irreducible
cubic \(\det T=0\), the determinant has multiplicity three along it,
and this actual nine-dimensional cross pencil has no rank-eight point.
\end{proposition}
\begin{proof}
For \(\operatorname{rank}T=2\), its columns give a two-dimensional
space of matrices in \(\langle A_0,A_1,A_2\rangle\).
This space contains an invertible matrix: the exact formula
\[
 \det(rA_0+sA_1+tA_2)=r^2t+rt^2+s^3
 \tag{PTC5}
\]
defines a smooth plane cubic, since its three partial derivatives
are \(t^2,s^2,r^2\). In particular it contains no projective line.
The right-coordinate transformations in (PTC4) can therefore
put the three column matrices into the form \(U,V,0\), with
\(U\) invertible. In these coordinates the cross matrix is
\[
                    \begin{pmatrix}0&U&V\\U&0&0\\V&0&0\end{pmatrix}.
 \tag{PTC6}
\]
Its kernel is exactly
\(\{(0,U^{-1}Vz,z):z\in k^3\}\), while its upper
six-dimensional minor has determinant \(\det(U)^2\).
Its rank is consequently six.

If \(\operatorname{rank}T=1\), then \(C(T)\) is the tensor
product of a three-dimensional matrix and a nonzero alternating
three-dimensional matrix of rank two; it has rank at most six.
The zero coefficient matrix gives rank zero. Matrices of rank three
give invertible \(C(T)\) by (PTC1), so rank eight never occurs.
Finally, the rank-two matrices form a dense irreducible orbit in
the usual determinant hypersurface; their adjugates are nonzero,
so this cubic is reduced and irreducible. Formula (PTC1) gives
the exact multiplicity three, with no inference from a restricted
Segre locus of functionals.
\end{proof}

\begin{remark}[Scope of the obstruction and the retained computation]
\label{rem:ss-product-one-axis-cross-cube-scope}
\status{Candidate scope; no global-image or cycle conclusion.}
Nonvanishing of an actual twisted cross determinant does not imply
a reduced corank-one determinant wall, even on a supersingular
elliptic product. The example is compatible with all previously
proved full product monodromy theorems. It concerns the dual cross
pencil in the displayed bases; no identification with a complete
Frobenius-cycle operator, no cycle stable-rank assertion, and no
indecomposable-surface factorization is inferred here.

The exact program
\path{code/monodromy/verify_ss_product_one_axis_determinant_cube.py}
checks the universal elliptic determinant, unit Schur complement,
exterior-square covariance and determinant character, an explicit
three-dimensional kernel and nonzero six-minor in (PTC6), and a
full nine-variable determinant at the nontrivial point \((a,b)=(0,0)\).
Its 36 terms independently agree with the binary cube of the
six-term three-by-three determinant. Uniformity in the twist is
proved by (PTC4) and the universal unit calculation, not inferred
from that single specialization. The first exploratory program
computed the same full polynomial and then encountered a SymPy
multivariate-gcd exception; its executed source, output and traceback
are retained unchanged in
\path{agents/monodromy/notes/CP35_product_determinant_gcd_failure}.
\end{remark}
